\documentclass[11pt,letterpaper]{article}

\usepackage[utf8]{inputenc}
\usepackage[T1]{fontenc}
\usepackage[spanish,es-noshorthands]{babel}
\usepackage{spectral}
\usepackage[margin=2.5cm,headheight=48pt,headsep=18pt]{geometry}
\usepackage{fancyhdr}
\usepackage{graphicx}
\usepackage{xcolor}
\usepackage{listings}
\usepackage{tikz}
\usepackage{booktabs}
\usepackage{enumitem}
\usepackage{float}
\usepackage[hidelinks]{hyperref}
\usetikzlibrary{positioning}

\newcommand{\curso}{CC4303-1 Redes}
\newcommand{\tituloinforme}{Actividad: Forwarding básico}

\renewcommand{\sectionmark}[1]{\markboth{#1}{}}
\setlength{\parskip}{0.6em}
\setlength{\parindent}{0pt}
\setlength{\emergencystretch}{3em}

% Cabecera y pie de página con el mismo estilo del informe de práctica.
\pagestyle{fancy}
\fancyhf{}
\fancyhead[L]{\small\leftmark}
\fancyfoot[L]{\small Informe \tituloinforme}
\fancyfoot[R]{\small CC4303-1}
\renewcommand{\headrulewidth}{0.4pt}
\renewcommand{\footrulewidth}{0.4pt}

\fancypagestyle{portada}{%
  \fancyhf{}
  \fancyhead[L]{\small\begin{tabular}[b]{@{}l@{}}
    Universidad de Chile\\
    Facultad de Ciencias Físicas y Matemáticas\\
    Departamento de Ciencias de la Computación\\
    \curso
  \end{tabular}}
  \fancyhead[R]{\includegraphics[height=1.3cm]{logo_fcfm.png}}
  \renewcommand{\headrulewidth}{0.4pt}
  \renewcommand{\footrulewidth}{0pt}
}

\definecolor{fondocodigo}{gray}{0.96}
\lstset{
  basicstyle=\ttfamily\footnotesize,
  backgroundcolor=\color{fondocodigo},
  frame=single,
  rulecolor=\color{gray!50},
  breaklines=true,
  columns=fullflexible,
  keepspaces=true,
  showstringspaces=false,
  literate={á}{{\'a}}1 {é}{{\'e}}1 {í}{{\'i}}1 {ó}{{\'o}}1 {ú}{{\'u}}1
           {Á}{{\'A}}1 {ñ}{{\~n}}1
}

% Diagramas de red: routers como nodos y enlaces como líneas.
\tikzset{
  router/.style={circle,draw,thick,minimum size=8mm,font=\small},
  default/.style={rectangle,draw,thick,rounded corners,minimum size=8mm,font=\small,fill=gray!10},
  link/.style={thick,<->,>=stealth},
}

\begin{document}

\begin{titlepage}
  \thispagestyle{portada}
  \centering
  \vspace*{5cm}
  {\LARGE Informe\par}
  \vspace{0.3cm}
  {\Large \tituloinforme\par}
  \vspace{0.4cm}
  {\large Juan Pablo Bustamante Flaño\par}
  \vfill
  Correo electrónico: juan.bustamante.f@ug.uchile.cl\\
  Curso: \curso{} --- Primavera 2026\\
  Profesora: Ivana Bachmann\\
  Módulo 4: Redes y Ruteo, semanas 11-12\\
  Fecha: 18/11/2026
  \vspace*{3cm}
\end{titlepage}

\section{Introducción}

Se simuló un mini-Internet en la capa de red: cada router es un proceso con un socket UDP
en (127.0.0.1, puerto) y una tabla de rutas en un archivo. Si un paquete va dirigido al
router, este imprime el mensaje; si no, lo reenvía según su tabla. El código final (con
TTL) está en \texttt{actividades/Forwarding}: \texttt{router.py}, \texttt{enviar.py},
\texttt{prueba\_router.py} y las tablas en \texttt{rutas/}. Las pruebas de la parte I se
hicieron con la versión intermedia sin TTL.

\section{Parte I: mini-Internet sin TTL}

\subsection{Tabla de rutas simulada versus real}

\textbf{Similitudes:} ambas asocian una red de destino (en notación CIDR) con un siguiente
salto, se consultan para cada paquete y descartan el paquete si ninguna ruta coincide.

\textbf{Diferencias:}
\begin{itemize}[nosep]
  \item En la simulación todos los routers comparten la IP \texttt{127.0.0.1} y las redes
        se distinguen por rangos de puertos; en una red real se distinguen por IP.
  \item La simulación usa la primera línea que coincide; una tabla real usa el prefijo
        más largo.
  \item Una tabla real incluye interfaz de salida y métrica, y se actualiza con protocolos
        de ruteo; la simulada es un archivo estático.
\end{itemize}

\subsection{Header y script de envío}

El header usa 4 bytes para la IP (uno por número), 2 bytes para el puerto y luego el
mensaje. \texttt{enviar.py} recibe la información por \texttt{sys.argv}: el destino final
en formato \texttt{IP;puerto} (en la parte II, \texttt{IP;puerto;TTL}), el mensaje y la
dirección a la que se envía el paquete:

\begin{lstlisting}
% python3 enviar.py "127.0.0.1;8885;10" "hola R5" 127.0.0.1 8881
\end{lstlisting}

El primer argumento va entre comillas porque en \texttt{bash} el \texttt{;} separa
comandos.

\subsection{Round-robin}

\texttt{check\_routes} revisa la tabla en orden. La primera línea que contiene al destino
define el \emph{área} de la red, dada por (red, puerto inicial, puerto final), y todas las
líneas con esa misma área son rutas alternativas. Un diccionario global guarda, para cada
área, el índice de la próxima ruta a usar, que avanza de forma circular con cada paquete.
Así cada área tiene su propio contador, para cualquier cantidad de rutas.

\subsection{Router default}

La ruta default debe ir en la \textbf{última línea}: la tabla se revisa en orden y la ruta
default coincide con cualquier destino, así que en cualquier otra posición taparía las
rutas que están debajo. Las tablas quedan así (ejemplo de R2):

\begin{lstlisting}
# rutas_R2_v5.txt
127.0.0.1/24 8881 8881 127.0.0.1 8881
127.0.0.1/24 8883 8885 127.0.0.1 8883
127.0.0.1/24 8883 8885 127.0.0.1 8884
0.0.0.0/0 0 65535 127.0.0.1 7000
\end{lstlisting}

Un paquete enviado desde R1 al puerto 9000, que no pertenece a la red, llegó al router
default:

\begin{lstlisting}
redirigiendo paquete 127.0.0.1;9000;fuera de la red con destino final ('127.0.0.1', 9000) desde ('127.0.0.1', 8881) hacia ('127.0.0.1', 7000)
\end{lstlisting}

\section{Pruebas sin TTL}

\subsection{Tabla de rutas mal configurada}

Con la tabla de R2 indicada en el enunciado, un paquete de R1 a R3 queda rebotando entre
R1 y R2 para siempre: en 2 segundos hubo casi 78\,000 reenvíos y el paquete nunca llegó a
R3. Un solo error de configuración genera un ciclo infinito que ocupa a los routers y al
enlace indefinidamente.

\subsection{Saltos con round-robin (5 routers)}

Se enviaron 12 paquetes de R1 a R5, de a uno:

\begin{lstlisting}
p01:  7 saltos  R1 -> R2 -> R3 -> R2 -> R4 -> R2 -> R3 -> R5
p02:  3 saltos  R1 -> R2 -> R4 -> R5
...   (el patrón se repite)
\end{lstlisting}

Los paquetes no siempre dan los mismos saltos: alternan entre 7 y 3, cuando el mínimo es 3.
El round-robin no considera la distancia, y por eso R3 a veces devuelve el paquete a R2.

\subsection{Saltos con round-robin (7 routers)}

Se agregaron R0 (8880), conectado a R1 y R2, y R6 (8886), conectado a R2 y R3
(figura~\ref{fig:v4}). Al intercalar paquetes de R1 a R5 y de R5 a R1, R2 cicló
correctamente dentro de cada área: R3 $\rightarrow$ R4 $\rightarrow$ R6 hacia R5, y
R0 $\rightarrow$ R1 hacia R1. Los caminos variaron entre 3 y 11 saltos, y solo 3 de 12
paquetes usaron el mínimo (3 saltos). Algunos pasaron varias veces por el mismo router.

\begin{figure}[H]
\centering
\begin{tikzpicture}[node distance=1.8cm]
  \node[router] (r1) {R1};
  \node[router,right=of r1] (r2) {R2};
  \node[router,right=of r2] (r3) {R3};
  \node[router,below=1.2cm of r2] (r4) {R4};
  \node[router,below=1.2cm of r3] (r5) {R5};
  \node[router,above=1.2cm of r1,xshift=1.4cm] (r0) {R0};
  \node[router,above=1.2cm of r2,xshift=1.4cm] (r6) {R6};
  \foreach \a/\b in {r1/r2,r2/r3,r2/r4,r3/r5,r4/r5,r0/r1,r0/r2,r6/r2,r6/r3}
    \draw[link] (\a) -- (\b);
\end{tikzpicture}
\caption{Red de 7 routers.}
\label{fig:v4}
\end{figure}

\begin{lstlisting}
# rutas_R0_v4.txt
127.0.0.1/24 8881 8886 127.0.0.1 8881
127.0.0.1/24 8881 8886 127.0.0.1 8882
# rutas_R1_v4.txt
127.0.0.1/24 8880 8880 127.0.0.1 8880
127.0.0.1/24 8882 8886 127.0.0.1 8882
127.0.0.1/24 8882 8886 127.0.0.1 8880
# rutas_R2_v4.txt
127.0.0.1/24 8880 8881 127.0.0.1 8880
127.0.0.1/24 8880 8881 127.0.0.1 8881
127.0.0.1/24 8883 8886 127.0.0.1 8883
127.0.0.1/24 8883 8886 127.0.0.1 8884
127.0.0.1/24 8883 8886 127.0.0.1 8886
# rutas_R3_v4.txt
127.0.0.1/24 8880 8882 127.0.0.1 8882
127.0.0.1/24 8880 8882 127.0.0.1 8886
127.0.0.1/24 8884 8885 127.0.0.1 8882
127.0.0.1/24 8884 8885 127.0.0.1 8885
127.0.0.1/24 8886 8886 127.0.0.1 8886
# rutas_R4_v4.txt
127.0.0.1/24 8880 8883 127.0.0.1 8882
127.0.0.1/24 8880 8883 127.0.0.1 8885
127.0.0.1/24 8885 8886 127.0.0.1 8882
127.0.0.1/24 8885 8886 127.0.0.1 8885
# rutas_R5_v4.txt
127.0.0.1/24 8880 8884 127.0.0.1 8884
127.0.0.1/24 8880 8884 127.0.0.1 8883
127.0.0.1/24 8886 8886 127.0.0.1 8883
# rutas_R6_v4.txt
127.0.0.1/24 8880 8882 127.0.0.1 8882
127.0.0.1/24 8880 8882 127.0.0.1 8883
127.0.0.1/24 8883 8885 127.0.0.1 8883
127.0.0.1/24 8883 8885 127.0.0.1 8882
\end{lstlisting}

\section{Parte II: mini-Internet con TTL}

El TTL se agregó como 1 byte después del puerto (header de 7 bytes). El router descarta
los paquetes que llegan con TTL 0 y, antes de reenviar uno, le resta 1. Con TTL 1 de R1 a
R3, R1 reenvió el paquete con TTL 0 y R2 lo descartó:

\begin{lstlisting}
Se recibió paquete 127.0.0.1;8883;0;hola con TTL 0
\end{lstlisting}

\texttt{create\_packet} retorna \texttt{bytes}, no \texttt{str} como supone el test del
enunciado, porque el header en bytes no siempre es texto UTF-8 válido. Para enviar un
archivo línea por línea:

\begin{lstlisting}
% python3 prueba_router.py "127.0.0.1;8885;10" 127.0.0.1 8881 < archivo.txt
\end{lstlisting}

\section{Pruebas con TTL}

\subsection{Tabla de rutas mal configurada}

Con TTL 10 el paquete sigue sin llegar a R3, pero el ciclo ya no es infinito: después de
10 reenvíos R1 lo recibe con TTL 0 y lo descarta. El TTL no corrige el error de
configuración, pero evita que el paquete quede circulando para siempre.

\subsection{Orden de los paquetes}

Se envió un archivo de 40 líneas de R1 a R5 con TTL 10. R5 recibió 37 líneas, en desorden:

\begin{lstlisting}
03 04 07 08 01 11 13 14 16 06 18 19 21 12 23 24 26 27 29 30 09 32 10 35 37 38 40
02 15 28 33 34 22 39 17 36 25
\end{lstlisting}

Por el round-robin, cada paquete toma un camino distinto, de entre 3 y 9 saltos, y los que
toman caminos cortos adelantan a los demás. Las líneas 05, 20 y 31 superaron los 9
reenvíos y se descartaron por TTL. Igual que IP, la red no garantiza ni el orden ni la
entrega de los paquetes.

\section{Casos borde}

\begin{itemize}[nosep]
  \item Como el TTL se revisa antes de ver el destino, un paquete que llega a su destino
        con TTL 0 también se descarta. Para llegar a $k$ saltos se necesita TTL $\geq k+1$.
  \item El router compara el destino con la IP recibida como argumento, así que debe
        iniciarse con \texttt{127.0.0.1} y no con \texttt{localhost}.
\end{itemize}

\section*{Declaración de uso de IA}

En este informe se utilizó el modelo de IA Claude exclusivamente como apoyo en la
revisión del código, recomendaciones de diseño, la redacción del informe y la
elaboración del diagrama en \LaTeX{}. Las decisiones de
diseño, la implementación del código, la ejecución de los experimentos y la obtención de
los datos fueron realizadas por el autor. Todo el contenido generado con apoyo de la
herramienta fue revisado y validado, y su responsabilidad es del autor.

\end{document}
